\section{Curto-circuito, seletividade e arco elétrico}

\begin{frame}{Corrente de curto-circuito presumida}
\small
\begin{columns}[T,onlytextwidth]
\column{0.50\textwidth}
\begin{align*}
I_{cc,3\phi}&\approx\frac{V_{LL}}{\sqrt3 Z_{eq}}\\
I_{cc,trafo}&\approx\frac{I_N}{Z_{pu}}
\end{align*}
\begin{itemize}
\item Fonte, transformador, cabos e conexões entram em $Z_{eq}$.
\item Calcular curto máximo e mínimo conforme finalidade.
\end{itemize}
\column{0.48\textwidth}
\begin{caixaEx}
Trafo $\SI{300}{kVA}$, $\SI{380}{V}$, $Z=5\%$:
\[
I_N=\frac{300000}{\sqrt3\cdot380}=\SI{456}{A}
\]
\[
I_{cc}\approx\frac{456}{0{,}05}=\SI{9.1}{kA}.
\]
É uma aproximação no secundário do trafo, antes da impedância dos cabos.
\end{caixaEx}
\end{columns}
\end{frame}

\begin{frame}{Capacidade de interrupção e suportabilidade}
\small
\begin{columns}[T,onlytextwidth]
\column{0.49\textwidth}
\begin{caixaAt}[title=Dispositivo]
A capacidade de interrupção do disjuntor/fusível deve ser compatível com a corrente de curto presumida no ponto.
\end{caixaAt}
\column{0.49\textwidth}
\begin{caixaNorma}[title=Conjunto]
Barramentos, cabos e quadros precisam suportar os efeitos térmicos/dinâmicos até a eliminação da falta.
\end{caixaNorma}
\end{columns}
\vspace{3mm}
\begin{caixaObs}
``Disjuntor 10 kA'' não garante que todo QD/QGBT tenha suportabilidade de 10 kA. O conjunto deve ser especificado e verificado como sistema.
\end{caixaObs}
\end{frame}

\begin{frame}{Seletividade: manter apenas a parte defeituosa fora}
\small
\begin{columns}[T,onlytextwidth]
\column{0.49\textwidth}
\begin{itemize}
\item Seletividade amperimétrica.
\item Seletividade temporal.
\item Seletividade energética.
\item Zonas/intertravamentos em proteção eletrônica.
\item Tabelas de coordenação do fabricante quando aplicáveis.
\end{itemize}
\column{0.49\textwidth}
\centering
\begin{tikzpicture}[scale=0.82,every node/.style={font=\scriptsize}]
\tikzset{d/.style={draw=corPrim,fill=corDef!7,rounded corners,minimum width=1.7cm,minimum height=0.55cm}}
\node[d] (dg) at (0,3) {DJ geral};
\node[d] (d1) at (-2.2,1.8) {DJ QD1};
\node[d] (d2) at (2.2,1.8) {DJ QD2};
\node[d] (c1) at (-2.2,0.6) {Carga 1};
\node[d] (c2) at (2.2,0.6) {Carga 2};
\draw[-{Stealth},thick] (dg)--(d1);\draw[-{Stealth},thick] (dg)--(d2);\draw[-{Stealth},thick] (d1)--(c1);\draw[-{Stealth},thick] (d2)--(c2);
\draw[corAt,thick,decorate,decoration={zigzag,segment length=4pt,amplitude=2pt}] (c1.south) -- ++(0,-0.65) node[below]{falta};
\node[corAt,align=center] at (0,0.15) {atua\\QD1};
\end{tikzpicture}
\end{columns}
\end{frame}

\begin{frame}{Curvas tempo--corrente}
\small
\begin{columns}[T,onlytextwidth]
\column{0.43\textwidth}
\begin{itemize}
\item Eixo horizontal: corrente ou múltiplo de $I_n$.
\item Eixo vertical: tempo de atuação.
\item Para seletividade, a curva a montante deve ficar adequadamente coordenada com a jusante no domínio de interesse.
\item Curto mínimo também importa: precisa haver atuação suficientemente rápida.
\end{itemize}
\column{0.55\textwidth}
\begin{tikzpicture}
\begin{axis}[width=6.3cm,height=4.2cm,xmode=log,ymode=log,xlabel={$I$},ylabel={$t$},xmin=1,xmax=100,ymin=0.01,ymax=1000,grid=both,legend pos=south west]
\addplot[corPrim,thick] coordinates {(1.1,800)(1.5,120)(2,25)(4,2)(8,0.05)};\addlegendentry{jusante}
\addplot[corAccent,thick,dashed] coordinates {(1.3,900)(2,180)(4,35)(8,3)(20,0.08)};\addlegendentry{montante}
\end{axis}
\end{tikzpicture}
\end{columns}
\end{frame}

\begin{frame}{Arco elétrico e energia incidente}
\small
\begin{columns}[T,onlytextwidth]
\column{0.52\textwidth}
\begin{caixaDef}
Arco elétrico é uma descarga sustentada através do ar/plasma. Pode gerar calor extremo, pressão, metal vaporizado, luz intensa e ruído.
\end{caixaDef}
\begin{itemize}
\item Risco não é avaliado apenas pela tensão nominal.
\item Tempo de atuação da proteção tem grande influência.
\item Distância de trabalho e configuração do equipamento importam.
\end{itemize}
\column{0.46\textwidth}
\begin{caixaNorma}[title=NR-10 futura --- vigência em 01/06/2027]
A redação aprovada pela Portaria MTE nº 737/2026, com vigência a partir de 01/06/2027, prevê, quando aplicável, que o projeto considere \textbf{estudo de energia incidente} para definir medidas de proteção contra os efeitos térmicos do arco elétrico. Até 31/05/2027, permanece vigente a redação anterior da NR-10.
\end{caixaNorma}
\end{columns}
\end{frame}

\begin{frame}{Como reduzir risco de arco elétrico}
\small
\begin{columns}[T,onlytextwidth]
\column{0.49\textwidth}
\begin{itemize}
\item Desenergizar sempre que possível.
\item Reduzir tempo de eliminação da falta.
\item Ajustes temporários de manutenção quando tecnicamente previstos.
\item Relés rápidos, ZSI, diferencial de barras, detecção de arco.
\item Operação remota e distância.
\end{itemize}
\column{0.49\textwidth}
\begin{itemize}
\item Compartimentação adequada.
\item Limitação de corrente quando aplicável.
\item Intertravamentos e portas adequadas.
\item Sinalização, análise de risco e procedimentos.
\item EPI baseado na avaliação e não como primeira medida de controle.
\end{itemize}
\end{columns}
\begin{caixaAt}
EPI é a última barreira. O projeto deve reduzir a energia e a probabilidade de exposição antes de depender do trabalhador.
\end{caixaAt}
\end{frame}

%=====================================================================
